\section{Compile-time ORM Architecture}
\label{sec:architecture}

The compile-time ORM architecture of the library is organised around a clear separation between the logical data model, the intermediate representation of queries, and the backend-specific materialisation layer. This layer carries the first principal scientific axis of the thesis: how the compiler can participate in modelling, constraining, and validating data access before the runtime phase begins. To achieve this, the architecture is built from several layers, each with its own responsibility and with strictly delimited interfaces.

The theoretical chapter showed that unification is possible at the level of the logical model and the query IR, but not by ignoring the differences between backend families. After the coroutine-theory chapter, the focus here deliberately returns to the compile-time core of the system: entity descriptions, the query IR, \code{connector\_trait<DB>}, and the capability mechanism. The library is therefore treated not merely as an idea, but as an organised system of types, contracts, and formal boundaries on admissible operations.

\subsection{Architectural goals and constraints}

The architecture was designed under several constraints acting simultaneously. First, the user-facing API must remain within idiomatic C++, without introducing a separate DSL and without relying on runtime reflection. Second, the query structure must be representable as a \code{constexpr} object so that a significant portion of validation can be performed by the compiler. Third, the backend layer must remain extensible without inheritance from a common virtual interface, because such an interface would hide important backend differences and move correctness checks toward runtime.

A fourth constraint follows from the multi-backend nature of the system itself. The library must support SQL and non-SQL connectors without false equivalence. This means that the architecture must simultaneously provide a shared API and express the limitations of a concrete backend in a formal and testable manner. The choice of a trait-based connector layer together with a capability mechanism follows naturally from this requirement.

\begin{figure}[H]
\centering
\begin{tikzpicture}[node distance=1.3cm and 1.4cm]
    \node[ormaccent] (app) {\textbf{Application layer}\\entity types and query composition};
    \node[ormblock, below=of app] (entity) {\textbf{Entity layer}\\\code{property}, \code{relationship}, \code{table\_name\_trait}};
    \node[ormblock, below=of entity] (ir) {\textbf{Query-IR layer}\\\code{field}, \code{Rule}, \code{select\_query}, \code{insert\_query}};
    \node[ormblock, below=of ir] (connector) {\textbf{Connector layer}\\\code{connector\_trait<DB>} and capability tags};
    \node[ormblock, below=of connector] (driver) {\textbf{Driver / runtime layer}\\SQL/BSON/Redis/Cypher/CQL materialisation};

    \draw[ormline] (app) -- (entity);
    \draw[ormline] (entity) -- (ir);
    \draw[ormline] (ir) -- (connector);
    \draw[ormline] (connector) -- (driver);
\end{tikzpicture}
\caption{Layered architecture of the library}
\label{fig:architecture_layers}
\end{figure}

Figure~\ref{fig:architecture_layers} shows that the library follows a consistent vertical organisation. Application code does not communicate directly with the driver; it works with C++ types and query builders. The driver, in turn, does not know domain semantics, but only the materialised form generated by the connector layer. This is precisely why the architecture can remain both strict and extensible.

\subsection{Layered organisation and responsibility separation}

The topmost layer is the application code, which defines entity types, constructs queries through \code{select}, \code{insert}, \code{update}, and \code{deleteq}, and submits them to \code{db<DB>}. The next layer is the logical model of the library --- the entity descriptions and the query IR. The third layer is the connector abstraction, where \code{connector\_trait<DB>} materialises the same logic into backend-specific syntax and behaviour. At the bottom stands the concrete driver, wire protocol, or mock implementation.

This architecture has an important practical effect. If the user replaces \code{SQLiteDB} with \code{MongoDB}, the query IR remains logically the same, but the connector interprets it as filter/projection execution rather than SQL rendering. This does not mean that every IR is admissible on every backend; it means that the \emph{form of abstraction} is shared, while admissibility is determined formally through the capability mechanism.

\begin{table}[H]
\centering
\small
\caption{Principal architectural layers and their responsibilities}
\label{tab:architecture_responsibilities}
\begin{tabularx}{\textwidth}{L{2.8cm}L{4.0cm}Y}
\toprule
\textbf{Layer} & \textbf{Main artefacts} & \textbf{Responsibility} \\
\midrule
Application & entity types, query-builder calls & describes domain logic and formulates operations in idiomatic C++ \\
Logical & \code{property}, \code{relationship}, \code{field}, \code{Rule} & models data and query structure independently of a concrete backend \\
Connector & \code{connector\_trait<DB>}, capability tags & checks admissibility and materialises the IR into backend behaviour \\
Runtime/driver & concrete connection handles and driver APIs & executes the materialised query and returns results \\
\bottomrule
\end{tabularx}
\end{table}

Table~\ref{tab:architecture_responsibilities} shows that the library does not rely on a single monolithic object that does everything. Instead, decisions are localised in separate layers. This matters both for extensibility and for academic analysis, because each claim can be traced to concrete types and modules.

\subsection{The entity layer as architectural foundation}

The entity layer is built around \code{property<T, Name>}, \code{relationship<...>}, and \code{table\_name\_trait<T>}. It plays a dual role. On the one hand, it represents the C++ domain model. On the other hand, it serves as a logical description of the persistent schema or structure. This means that the architecture does not construct a separate runtime metadata registry. Instead, the type system contains enough information to derive field names, value types, relationships, and the target storage container.

Especially important is the fact that the entity layer is independent of the database family. Neither \code{property} nor \code{relationship} is SQL-specific. They describe logical attributes and logical relations. This creates the foundation for comparable behaviour across relational and non-relational backends. At the architectural level, this is exactly what makes it possible for the thesis to argue that the C++ model is the primary carrier of the logical schema.

\subsection{Query IR as the central architectural component}

The query IR is the central architectural component. It combines field references, predicates, placeholders, and additional clauses such as \code{order\_by}, \code{group\_by}, and \code{limit}. Instead of working with textual queries, the library constructs typed objects whose template parameters encode the logical structure of the operation.

From the architectural point of view, this has several consequences. First, structural validation occurs early. Second, the query IR can be analysed by the connector layer without losing semantic information. Third, prepared queries can preserve not an SQL string, but the intermediate representation itself, which is more natural for a library oriented around the C++ type system.

Fourth, the query IR is the single zone in which entity description and later backend materialisation meet. In that sense it is the architectural hinge of the whole library: high enough to remain storage-agnostic, yet concrete enough to be executed.

\subsection{The connector layer and the formal contract}

Formally, the connector layer is implemented through specialisation of \code{connector\_trait<DB>}. This approach is analogous to standard trait templates such as \code{std::iterator\_traits}. The type \code{DB} may be a simple wrapper over a connection handle, while the entire ORM logic remains in the trait specialisation. In this way, the library does not require the user or the connector author to inherit from a common base class.

\code{connector\_trait<DB>} defines at least three central aspects: \code{wire\_type}, \code{cursor\_type}, and \code{execute(...)}. The specialisation may also declare capability tags such as \code{supports\_joins}, \code{supports\_transactions}, \code{supports\_aggregation}, or \code{supports\_concurrent\_execute}. The presence or absence of these pseudo-markers is sufficient for the upper layers to impose compile-time restrictions.

\begin{lstlisting}[language=C++,caption={Excerpt from the formal connector contract},label={lst:connector_contract}]
template <typename DB>
concept is_connector = requires {
    typename connector_trait<DB>::template wire_type<int>::type;
    typename connector_trait<DB>::cursor_type;
};

template <typename DB, typename Cap>
inline constexpr bool has_capability =
    detail::capability_check<DB, Cap>::value;
\end{lstlisting}

Listing~\ref{lst:connector_contract} shows that the contract is structural rather than inheritance-based. This is a decisive architectural choice. If a virtual interface existed with runtime dispatch, some of the type-checkable information would be lost. The library instead prefers a compile-time structure in which the requirements remain visible both to the compiler and to the connector author.

\begin{figure}[H]
\centering
\begin{tikzpicture}[node distance=1.35cm and 1.25cm]
    \node[ormaccent] (query) {\textbf{Typed query}\\join, transaction, aggregation intent};
    \node[ormblock, below=of query] (db) {\code{db<DB>} facade};
    \node[ormblock, below left=of db] (capok) {Capability present\\query remains admissible};
    \node[ormdanger, below right=of db] (capfail) {Capability missing\\compile-time rejection};
    \node[ormblock, below=of capok] (exec) {\code{connector\_trait<DB>::execute}};

    \draw[ormline] (query) -- (db);
    \draw[ormline] (db) -- (capok);
    \draw[ormline] (db) -- (capfail);
    \draw[ormline] (capok) -- (exec);
\end{tikzpicture}
\caption{Capability gating between the query IR and the connector layer}
\label{fig:capability_gating}
\end{figure}

Figure~\ref{fig:capability_gating} is important because it shows where the architecture decides admissibility. This is not the driver and not a runtime log. The decision is taken in the public API layer, which already has access both to the query IR and to the trait-level information about the connector.

\begin{table}[H]
\centering
\small
\caption{Role of the principal capability tags in the architecture}
\label{tab:capability_tags}
\begin{tabularx}{\textwidth}{L{3.6cm}Y}
\toprule
\textbf{Capability} & \textbf{Architectural meaning} \\
\midrule
\code{supports\_joins} & admits join-based query forms for the selected backend \\
\code{supports\_transactions} & allows use of \code{transaction\_guard} together with \code{begin}/\code{commit}/\code{rollback} hooks \\
\code{supports\_aggregation} & marks support for aggregation operations in the logical query layer \\
\code{supports\_embedding} & makes native support for embedded structures explicit \\
\code{supports\_upsert} & marks backends for which upsert belongs to the natural execution contract \\
\code{supports\_bulk\_insert} & denotes the existence of a more efficient bulk-write path \\
\bottomrule
\end{tabularx}
\end{table}

\subsection{The user-facing db<DB> facade}

The class \code{db<DB>} is the public facade layer. It encapsulates a reference to a concrete backend object and provides a unified execution API. In practice, this layer has two roles. The first is convenience for the user: the library is used through the same interface regardless of backend. The second is architectural filtering: this is precisely where compile-time checks for capabilities and backend compatibility are performed.

Consequently, \code{db<DB>} is not merely a thin wrapper over \code{execute}. It is the place where the library formalises the boundary between admissible and inadmissible behaviour for the selected backend. It is exactly this layer that enables error messages such as ``this connector does not support \code{JOIN} operations'' instead of allowing such errors to surface later as ambiguous runtime behaviour.

\begin{lstlisting}[language=C++,caption={Excerpt from the \code{db<DB>} facade layer},label={lst:db_facade}]
template <typename DB>
    requires is_connector<DB>
class db
{
public:
    template <typename Query>
    auto operator<<(Query q)
    {
        if constexpr (is_select_query<Query>)
        {
            if constexpr (decltype(q.join_clauses())::size > 0)
            {
                static_assert(has_capability<DB, cap::supports_joins>,
                    "orm::db: this connector does not support JOIN operations.");
            }
        }
        return connector_trait<DB>::execute(*conn_, std::move(q));
    }
};
\end{lstlisting}

Listing~\ref{lst:db_facade} shows why \code{db<DB>} is more than a convenient facade. It is the formal location where query structure and connector capability information meet. This is where the compile-time architecture becomes a real limitation on API use.

\subsection{Migrations as an architectural extension}

The architecture also includes a prototype migration layer through \code{migrate<DB>}, \code{live\_schema}, \code{entity\_meta}, and DDL operations such as \code{create\_table\_op}, \code{add\_column\_op}, \code{drop\_column\_op}, and \code{alter\_column\_type\_op}. This layer is particularly important from the standpoint of the thesis because it makes explicit the relation between the C++ model and the real storage schema.

The migration layer is not uniformly complete for every backend, but architecturally it demonstrates that the library is not limited to query execution. With a sufficiently rich \code{connector\_trait<DB>} specialisation, entity description can also participate in schema evolution. This changes the very meaning of the ORM layer: it becomes not only an interface to data, but a potential source of schema-aware tooling.

\subsection{Implementation of the principal subsystems}

Once the architectural model has been outlined, the next question is how it is materialised in concrete C++ constructs and how the subsystems cooperate inside the repository. This section restores the missing engineering perspective: not merely which layers exist, but how \code{where} clauses are accumulated, how a query object is preserved for reuse, how \code{migrate<DB>} turns the difference between entity description and \code{live\_schema} into DDL operations, and how the helper layer enforces rules for thread safety and transaction handling.

This section also fixes the working notation without reintroducing a standalone glossary chapter. Here, \code{DB} denotes a generic backend type, \code{db<DB>} denotes the public execution facade, \code{connector\_trait<DB>} is the formal connector contract, \code{field<...>} is the compile-time representation of a field, \code{Placeholder<T>} is an anonymous runtime parameter, and \code{orm::ph<T, \_N>} is an indexed placeholder. The repository context for this mechanics is concentrated mainly in \path{lib/include/ORM/entity/}, \path{lib/include/ORM/query/}, \path{lib/include/ORM/connector/}, \path{lib/include/ORM/db/migration/}, and in the verification layer under \path{tests/unit/} and \path{tests/integration/}.

\subsubsection{Mini case study with User and Post}

To keep the exposition directly tied to the real code, this section uses the \code{User} and \code{Post} types from \path{tests/integration/test_mockdb.cpp} as a supporting mini case study. They are especially useful because they combine scalar fields, explicit naming of logical containers, and a join scenario. The same example is later reused for placeholder binding, prepared-query behaviour, and SQL rendering assertions.

\begin{lstlisting}[language=C++,caption={Entity types from \protect\path{tests/integration/test_mockdb.cpp}},label={lst:mockdb_entities}]
struct Post
{
    orm::property<int, "id">             id;
    orm::property<int, "author_id">      author_id;
    orm::property<std::u8string, "body"> body;
};

struct User
{
    orm::property<int, "id">            id;
    orm::property<std::u8string, "name"> name;
    orm::property<double, "score">      score;
};
\end{lstlisting}

Listing~\ref{lst:mockdb_entities} shows not only how the entity model is described, but also how it participates directly in the tests. \code{table\_name\_trait<User>} and \code{table\_name\_trait<Post>} fix the logical containers \code{"users"} and \code{"posts"}, while member pointers such as \code{\&User::id} and \code{\&Post::author\_id} later participate in the query builder. This is a direct illustration of one of the central motifs of the thesis: the model, the query form, and the execution tests all use the same compile-time artefacts.

\subsubsection{Scalar properties and naming of containers}

The implementation of \code{property<T, Name>} binds a value type to a compile-time name. The function \code{column\_name()} gives unambiguous access to the symbolic field name and is used by connectors when rendering toward SQL columns, document fields, key-space segments, or other backend-specific notions. The construction is minimalistic but sufficient because it avoids any external runtime metadata registry.

Naming of the logical storage container is implemented through \code{table\_name\_trait<T>}. Although the name is inherited from ORM tradition, architecturally it plays a more general role: in SQLite, PostgreSQL, and MySQL it means a table, in MongoDB a collection, in Redis a key prefix, and in Neo4j a label. In this way the library uses a unified logical interface for addressing physically different structures.

The practical consequence matters: a connector author does not derive schema information from an external registry, but from the type system. This reduces boilerplate and makes naming mistakes easier to localise both at compile time and in unit tests.

\subsubsection{Field wrappers, rules, and placeholders}

The field in a query expression is represented by \code{field<\&T::m>}. This is a \code{constexpr} wrapper over a member pointer that carries enough information about the container type, the value type, and the column name. Operators over \code{field} construct typed \code{Rule} objects. Expressions such as \code{field<\&User::id> == Placeholder<int>\{\}} are therefore not strings, but compile-time descriptions of predicates.

The placeholder mechanism has two forms. \code{Placeholder<T>} models anonymous parameters that are consumed positionally by the arguments of \code{execute()}. \code{orm::ph<T, std::placeholders::\_N>} models indexed parameters that allow one runtime value to be reused in more than one position. This is especially important for backends such as SQLite, PostgreSQL, and MySQL, where the placeholder slot is a distinct part of the materialisation logic.

The tests \code{IndexedPlaceholderRendersQuestionMarkN}, \code{IndexedPlaceholderReusedSameArgInSQL}, and \code{TwoDistinctIndexedPlaceholdersRenderCorrectSlots} in \path{tests/integration/test_mockdb.cpp} show that this subsystem is not a theoretical addition, but a real mechanism for expressing more precise parameter-binding scenarios.

\subsubsection{Construction of query objects}

The factory functions \code{select}, \code{insert}, \code{update}, and \code{deleteq} create query objects upon which \code{where}, \code{join}, \code{order\_by}, \code{group\_by}, and \code{limit} are layered in sequence. Crucially, each of these operations does not modify a string, but returns a new typed object with extended internal state. As a result, the query builder remains fluent in appearance yet static in nature.

This allows the tests to validate not only the final rendering result, but also the form of the intermediate representation itself --- for example the number of \code{where} clauses, the response-tuple type, or the correctness of \code{GROUP BY} and \code{ORDER BY} composition. In \path{tests/integration/test_mockdb.cpp} this is visible in \code{SelectWithInnerJoin}, \code{SelectWithOrderByDesc}, \code{GroupByMultipleFields}, \code{SelectWithLimit}, and numerous other scenarios that validate different phases of query composition.

One key implementation detail is that the builder chain remains \code{constexpr}-compatible. The internal helper \code{orm::detail::tuple\_cat} uses pack expansion through index sequences rather than reliance on runtime concatenation. The practical consequence is that a significant part of the query shape can be built and validated already during compilation.

\begin{lstlisting}[language=C++,caption={Query object declared as a compile-time constant},label={lst:constexpr_query}]
using namespace orm::literals;
static constexpr auto get_active =
    orm::select(orm::field<&User::id>, orm::field<&User::name>)
        .where(orm::field<&User::score> > 0.0)
        .where(orm::field<&User::id> == orm::Placeholder<int>{})
        .limit(10_per_page & 2_page);
\end{lstlisting}

The operator \code{\&} combines the \code{\_per\_page} and \code{\_page} helpers into a \code{Pagification} object whose structure is known at compile time. The header \path{lib/include/ORM/query/limits.hpp} also supports a symmetric \code{*} form, but the use of \code{\&} in Listing~\ref{lst:constexpr_query} is more readable for the query style of the library. In this way the query can be declared as \code{static constexpr} and reused without reconstructing the builder form at every call.

\begin{figure}[H]
\centering
\begin{tikzpicture}[node distance=1.35cm and 1.2cm]
    \node[ormaccent] (api) {\textbf{Fluent API}\\\code{select}/\code{where}/\code{join}/\code{limit}};
    \node[ormblock, below=of api] (typed) {New typed\\query state at each step};
    \node[ormblock, below left=of typed] (shape) {Compile-time\\query structure};
    \node[ormblock, below right=of typed] (params) {Runtime\\parameter values};
    \node[ormblock, below=of $(shape)!0.5!(params)$] (exec) {\code{connector\_trait<DB>::execute}};

    \draw[ormline] (api) -- (typed);
    \draw[ormline] (typed) -- (shape);
    \draw[ormline] (typed) -- (params);
    \draw[ormline] (shape) -- (exec);
    \draw[ormline] (params) -- (exec);
\end{tikzpicture}
\caption{Interaction between query composition, compile-time structure, and runtime parameters}
\label{fig:query_composition_execution}
\end{figure}

Figure~\ref{fig:query_composition_execution} shows that the implementation operates in two modes at once: the structural form of the query is fixed in the type, while the concrete values are supplied separately to \code{execute(...)}. This is exactly what makes the coexistence of compile-time validation and runtime parameterisation possible.

\subsubsection{Execution, results, and field access}

At execution time, \code{db<DB>} forwards the query IR to \code{connector\_trait<DB>::execute}. The result is returned as \code{orm::result<Row, FieldTuple>}. In this construction, \code{Row} is the materialised tuple type, while \code{FieldTuple} preserves the projection metadata for the selected fields. This solution allows both iteration over materialised rows and access through \code{get\_field<\&T::m>} or positional access through \code{get<I>}. The implementation therefore remains type-safe even after query execution.

The tests \code{SelectResultRowTypeIsTuple}, \code{SelectMultiFieldRowType}, \code{SelectGetByMemberPointer}, and \code{SelectGetByIndex} show that the result layer is not merely a container of anonymous rows, but a structured abstraction with type-traceable projection information. This matters for the thesis because it demonstrates that type discipline does not end with the query builder, but continues into result handling.

\subsubsection{Prepared queries as a reusable execution artefact}

\code{prepared\_query<DB, Query>} extends the execution model. Instead of caching an SQL string, it stores a reference to the backend object and the query IR itself. This is a natural continuation of the compile-time approach: the library treats the query object as the primary artefact rather than as a transient step on the way to a textual result.

\begin{lstlisting}[language=C++,caption={Excerpt from \code{prepared\_query<DB, Query>}},label={lst:prepared_query}]
template <typename DB, typename Query>
class prepared_query
{
public:
    prepared_query(DB& conn, Query q) : conn_(conn), query_(std::move(q)) {}

    template <typename... Params>
    auto execute(Params&&... params) const
    {
        return connector_trait<DB>::execute(
            conn_, query_, std::forward<Params>(params)...);
    }

    [[nodiscard]] const Query& query() const noexcept { return query_; }
};
\end{lstlisting}

Listing~\ref{lst:prepared_query} reveals an important engineering choice: reuse is organised around the IR rather than around textual materialisation. This allows the same logical query to be executed repeatedly with different parameters without reconstructing the builder form. This behaviour is validated exactly by \code{PrepareReturnsExecutableObject}, \code{PrepareWithParamExecutesCorrectly}, \code{PrepareCanBeCalledMultipleTimesWithDifferentParams}, and \code{PrepareExposesQueryIR}.

\begin{figure}[H]
\centering
\begin{tikzpicture}[node distance=1.35cm and 1.2cm]
    \node[ormaccent] (build) {Construction of query IR};
    \node[ormblock, right=of build] (prepare) {\code{db.prepare(query)}};
    \node[ormblock, right=of prepare] (store) {Stored IR +\\reference to \code{DB}};
    \node[ormblock, below=of store] (exec1) {\code{execute(1)}};
    \node[ormblock, left=of exec1] (exec2) {\code{execute(99)}};
    \node[ormblock, right=of exec1] (exec3) {\code{execute(7)}};

    \draw[ormline] (build) -- (prepare);
    \draw[ormline] (prepare) -- (store);
    \draw[ormline] (store) -- (exec1);
    \draw[ormline] (store) -- (exec2);
    \draw[ormline] (store) -- (exec3);
\end{tikzpicture}
\caption{Lifecycle of \code{prepared\_query}: one IR, multiple executions}
\label{fig:prepared_query_lifecycle}
\end{figure}

\subsubsection{Migration layer and DDL generation}

The implementation of \code{migrate<DB>} compares the schema described by entity types against \code{live\_schema}. Whenever a difference exists, DDL operations are created and forwarded to \code{connector\_trait<DB>::ddl\_for(...)}. This is a good example of clean separation of responsibilities: the diff logic is shared, whereas the concrete DDL syntax remains backend-specific.

\begin{lstlisting}[language=C++,caption={Core of the migration diff pipeline},label={lst:migration_diff}]
[[nodiscard]] std::vector<ddl_op> diff(
    const live_schema& live,
    std::span<const entity_meta> entities) const
{
    std::vector<ddl_op> ops;

    for (const auto& entity : entities)
    {
        auto live_it = live.find(entity.table_name);
        if (live_it == live.end())
        {
            ops.push_back(create_table_op{entity.table_name, entity.columns});
            continue;
        }

        for (const auto& [col, type] : entity.columns)
        {
            if (live_it->second.find(col) == live_it->second.end())
                ops.push_back(add_column_op{entity.table_name, col, type});
        }
    }
    return ops;
}
\end{lstlisting}

The distinction between a \emph{minimally working connector} and a \emph{full connector contract} becomes especially clear here. A backend may execute queries without supporting DDL generation. But if it is expected to participate in schema-aware tooling, its contract must be extended with \code{ddl\_for(...)} hooks and consistent migration verification.

\begin{figure}[H]
\centering
\begin{tikzpicture}[node distance=1.35cm and 1.2cm]
    \node[ormaccent] (entity) {Entity metadata\\\code{entity\_meta}};
    \node[ormblock, right=of entity] (live) {Runtime schema\\\code{live\_schema}};
    \node[ormblock, below=of $(entity)!0.5!(live)$] (diff) {\code{migrate<DB>::diff}};
    \node[ormblock, below left=of diff] (ops) {DDL operations\\\code{create/add/drop/alter}};
    \node[ormblock, below right=of diff] (ddl) {\code{generate\_ddl}};
    \node[ormblock, below=of $(ops)!0.5!(ddl)$] (run) {dry-run or execution path};

    \draw[ormline] (entity) -- (diff);
    \draw[ormline] (live) -- (diff);
    \draw[ormline] (diff) -- (ops);
    \draw[ormline] (diff) -- (ddl);
    \draw[ormline] (ops) -- (run);
    \draw[ormline] (ddl) -- (run);
\end{tikzpicture}
\caption{Flow of the schema-migration subsystem}
\label{fig:migration_pipeline}
\end{figure}

The tests in \path{tests/unit/test_schema_migration.cpp} cover table creation, column addition and removal, dry-run completion codes, DDL generation, and the state-machine behaviour of \code{run()}. This is strong evidence that the migration layer is not merely a concept, but a genuinely traceable pipeline.

\subsubsection{Thread safety and transaction helpers}

The library contains a dedicated layer for \code{connection\_pool<DB, N>}, \code{thread\_local\_db<DB>}, and \code{transaction\_guard<DB>}. These components show that the implementation does not stop at one-off query execution, but also considers operational scenarios involving connection reuse, per-thread isolation, and rollback discipline. The important architectural choice is that access to these mechanisms is likewise capability-based. For example, \code{connection\_pool} requires \code{supports\_concurrent\_execute}, whereas \code{transaction\_guard} requires \code{supports\_transactions}.

\begin{lstlisting}[language=C++,caption={Excerpt from the thread-safety helper layer},label={lst:thread_safety_helpers}]
template <typename DB, std::size_t N>
class connection_pool
{
    static_assert(
        requires { typename connector_trait<DB>::supports_concurrent_execute; },
        "connection_pool requires "
        "connector_trait<DB>::supports_concurrent_execute. "
        "Add 'using supports_concurrent_execute = void;' "
        "to connector_trait<DB>.");

public:
    [[nodiscard]] connection_guard<DB> acquire();
};

template <typename DB>
class transaction_guard
{
    static_assert(
        requires { typename connector_trait<DB>::supports_transactions; },
        "transaction_guard requires "
        "connector_trait<DB>::supports_transactions. "
        "Add 'using supports_transactions = void;' "
        "to connector_trait<DB>.");
};
\end{lstlisting}

In this way the compile-time contract is not restricted to query shape alone, but extends to the lifecycle policies of the connector. This is important for the thesis argument that the library is not merely a query builder, but infrastructure for systematic and safe work with backend resources.

\begin{table}[H]
\centering
\small
\caption{Principal subsystems and their test support}
\label{tab:core_subsystem_tests}
\begin{tabularx}{\textwidth}{L{3.1cm}L{4.1cm}Y}
\toprule
\textbf{Subsystem} & \textbf{Principal files} & \textbf{Test support} \\
\midrule
Entity modelling & \path{ORM/entity/property.hpp}, \path{ORM/entity/relationship.hpp} & \path{tests/unit/test_property.cpp}, \path{tests/unit/test_relationship.cpp} \\
Query composition & \path{ORM/query/select.hpp} and the related query headers & \path{tests/unit/test_query.cpp}, \path{tests/integration/test_mockdb.cpp} \\
Prepared queries & \path{ORM/connector/prepared_query.hpp} & the \code{Prepare*} scenarios in \path{tests/integration/test_mockdb.cpp} \\
Schema migration & \path{ORM/db/migration/migration.hpp} & \path{tests/unit/test_schema_migration.cpp} \\
Thread safety & \path{ORM/db/connectors/ThreadSafety/thread_safety.hpp} & \path{tests/unit/test_thread_safety.cpp} \\
\bottomrule
\end{tabularx}
\end{table}

\subsubsection{Implementation conclusion}

The principal subsystems have direct support in both code and tests. \path{tests/unit/test_property.cpp} and \path{tests/unit/test_relationship.cpp} validate the entity descriptions. \path{tests/unit/test_query.cpp} proves that the query builder genuinely constructs correct compile-time structures. \path{tests/integration/test_mockdb.cpp} and \path{tests/integration/test_sqlite.cpp} validate execution, prepared-query scenarios, and result handling. \path{tests/unit/test_schema_migration.cpp} covers migration diff logic, while \path{tests/unit/test_thread_safety.cpp} checks the lifecycle layer for concurrent access.

Consequently, implementation of the principal subsystems is not merely a design intention, but an architectural decision backed by systematic verification. This is the most important conclusion of the present section: the compile-time ORM architecture in the repository does not remain at the level of abstraction, but materialises into working subsystems with clearly separated responsibilities and traceable repository evidence.

\subsection{Architectural invariants}

On the basis of the analysed layers, several invariants can be formulated. First, application code should never depend directly on driver APIs. Second, query logic is formulated in a typed IR rather than in backend-native strings. Third, extension toward a new backend proceeds through trait specialisation and capability declarations rather than through changes to the query API itself. Fourth, if an operation is inadmissible for a backend, this should become visible as early as possible --- ideally at compile time.

These invariants make the architecture both engineering-sound and academically defensible. They allow every future extension to be judged not only by whether it ``works'', but by whether it preserves the already formulated layers and boundaries.

\subsection{Architectural conclusion}

In summary, the architecture of the library separates the problem into three levels: logical model, typed query IR, and backend materialisation. This separation enables strict type discipline, formal extensibility, and controlled behaviour across distinct storage models. The next chapter shifts the focus to the second architectural axis --- asynchronous execution --- where it becomes clear how the same compile-time core is connected to coroutine lifecycle, thread-pool offload, and native non-blocking integration.
